\section{Training budgets and model capacity}
\label{app:fairness}

A comparison between learned methods is only informative if the methods were
given comparable resources. Table~\ref{tab:fairness} lists what each one
actually received. Every learned method trains on the same curriculum, draws
instances from the same generator, and is evaluated zero-shot on the same
twelve held-out configurations under the same protocol.

\begin{table}[htbp]
\centering
\small
\caption{What each learned method was given. \emph{Instances seen} counts
training instances drawn, not gradient steps: our method draws $P{=}10$
rollouts per instance, so it performs more rollouts than instances. AM and
POMO use the RL4CO defaults for their architectures; we did not reduce them.}
\label{tab:fairness}
\begin{tabular}{@{}lrrrr@{}}
\toprule
 & Parameters & Encoder & Instances seen & Learning rate \\
\midrule
AM~\cite{kool2018attention}   & 694{,}528   & Transformer, 3 layers & 512{,}000 & $1\times10^{-4}$ \\
POMO~\cite{kwon2020pomo}      & 1{,}289{,}344 & Transformer, 6 layers & 128{,}000 & $1\times10^{-4}$ \\
Sequential                    & 1{,}291{,}394 & GNN, 3 layers         & 15{,}000  & $5\times10^{-4}$ \\
\textbf{Ours}                 & 1{,}307{,}778 & GNN, 3 layers         & \textbf{6{,}000} & $5\times10^{-4}$ \\
\bottomrule
\end{tabular}
\end{table}

Two things are worth drawing out. Capacity is within 2\% across our method,
Sequential and POMO; AM is smaller only because RL4CO's attention model
defaults to three encoder layers, and we left that default alone rather than
enlarging a baseline to match ours. Training budget runs the other way: AM
draws roughly 85 times as many instances as we do and POMO roughly 21 times,
because both use the batch sizes their own implementations ship with. The
proposed method is therefore not the one that received the most training, and
the comparison does not rest on having trained it harder.

No method here received hyperparameter tuning. The values in
Table~\ref{tab:hyperparameters} and the RL4CO defaults were used as they
stand, and the only parameter this paper sweeps is the inference-time search
budget $K$ (Section~\ref{sec:budget}), which no baseline uses. A tuned
baseline could close some of the gap; we report the untuned comparison
because tuning every method to the same standard was beyond what we could do,
and because tuning only our own would have been worse.

\section{What the exact solver actually returns}
\label{app:scip}

The SCIP column of Table~\ref{tab:main_results} is not an optimum except at
the two smallest configurations. SCIP is given 600\,s per instance and reports
its incumbent when the budget runs out, so the comparison is against a
time-limited incumbent almost everywhere. Table~\ref{tab:scip} records where
the line falls.

\begin{table}[htbp]
\centering
\small
\caption{SCIP under a 600\,s limit on the ten reported instances per
configuration. \emph{Proved} counts instances closed to optimality. The final
column is the mean relative gap on the instances SCIP did not close; entries
marked $\gg$ are configurations where no useful dual bound was found at all,
leaving the reported gap effectively unbounded.}
\label{tab:scip}
\begin{tabular}{@{}lrrr@{}}
\toprule
Config & Proved & Mean time & Gap if unproved \\
\midrule
(5,5,5)     & 10/10 & 5.2\,s   & --- \\
(5,7,5)     & 10/10 & 3.4\,s   & --- \\
(10,15,5)   & 6/10  & 279.4\,s & 20.2\% \\
\midrule
(15,15,5)   & 0/10  & 600.0\,s & $\gg$ \\
(15,20,5)   & 0/10  & 600.0\,s & 71.1\% \\
(20,30,5)   & 1/10  & 589.5\,s & 27.6\% \\
\midrule
(30,30,10)  & 0/10  & 600.0\,s & $\gg$ \\
(30,40,10)  & 0/10  & 600.0\,s & $\gg$ \\
(40,50,10)  & 0/10  & 600.4\,s & $\gg$ \\
\midrule
(50,50,15)  & 0/10  & 601.4\,s & $\gg$ \\
(50,70,15)  & 0/10  & 601.8\,s & $\gg$ \\
(70,100,15) & 0/10  & 604.6\,s & $\gg$ \\
\bottomrule
\end{tabular}
\end{table}

Two readings follow. At the Small tier the comparison is against a genuine
optimum, and SCIP wins there, as it should: 26 of the 30 instances in that
tier are closed. From (15,15,5) onward almost nothing is closed, the solver
spends its full budget on every instance, and on most configurations it
returns no dual bound worth reporting, so its column should be read as a
strong incumbent rather than a bound on what is achievable. That is also why
SCIP's objective at the two largest configurations, 0.992 and 0.993, is close
to leaving the entire target set intact: with 600\,s it does not reach a
competitive solution at that size at all.

The practical reading is that an exact solver remains the right tool where it
finishes. The case for a learned policy is confined to the region where it
does not, and Table~\ref{tab:scip} marks where that region begins.

\section{Evaluation traces}
\label{app:traces}

Table~\ref{tab:traces} gives every validation evaluation recorded during
training, for all ten seeds. It is included because the checkpoint reported in
Section~\ref{sec:results} is selected by best validation performance rather
than taken at the final step, and a reader is entitled to see what that
selection was choosing among. Validation instances are drawn at the training
size range under a seed used nowhere else, so none of the twelve reported
configurations appears here.

\begin{table*}[htbp]
\centering
\scriptsize
\caption{Validation objective at every evaluation during training, ten seeds,
$K{=}10$. Lower is better. Training is non-monotone: the best value for a run
is often far from its last, and three runs were still improving when the
schedule ended.}
\label{tab:traces}
\begin{tabular}{@{}lcccccccccc@{}}
\toprule
Step & s1 & s2 & s3 & s4 & s5 & s6 & s7 & s8 & s9 & s10 \\
\midrule
20 & 0.1689 & 0.1926 & 0.2001 & 0.5163 & 0.2148 & 0.2117 & 0.5393 & 0.2005 & 0.2642 & 0.1936 \\
40 & 0.1962 & 0.4990 & 0.2070 & 0.2163 & 0.1651 & 0.1941 & 0.2239 & 0.1731 & 0.1827 & 0.1533 \\
60 & 0.1859 & 0.2108 & 0.1643 & 0.2063 & 0.2744 & 0.1926 & 0.3947 & 0.2376 & 0.5401 & 0.3128 \\
80 & 0.2233 & 0.1777 & 0.3230 & 0.2005 & 0.2219 & 0.1820 & 0.2461 & 0.2258 & 0.1693 & 0.2685 \\
100 & 0.4949 & 0.3198 & 0.3205 & 0.1605 & 0.1942 & 0.2343 & 0.4087 & 0.3273 & 0.1623 & 0.2495 \\
120 & 0.2907 & 0.3715 & 0.3540 & 0.1745 & 0.2057 & 0.2122 & 0.2220 & 0.1597 & 0.2053 & 0.1866 \\
140 & 0.3997 & 0.1589 & 0.3254 & 0.5385 & 0.4319 & 0.3201 & 0.1535 & 0.1841 & 0.2330 & 0.2476 \\
160 & 0.1985 & 0.1603 & 0.2009 & 0.1729 & 0.2949 & 0.1681 & 0.1808 & 0.1672 & 0.2106 & 0.3447 \\
180 & 0.2316 & 0.2267 & 0.2281 & 0.1686 & 0.2193 & 0.4487 & 0.1500 & 0.3925 & 0.1797 & 0.3311 \\
200 & 0.2345 & 0.2300 & 0.2257 & 0.3212 & 0.1657 & 0.5255 & 0.3614 & 0.2551 & 0.2838 & 0.3320 \\
220 & 0.2873 & 0.2401 & 0.2989 & 0.1683 & 0.1836 & 0.4038 & 0.4939 & 0.1543 & 0.1565 & 0.3077 \\
240 & 0.1620 & 0.2314 & 0.2977 & 0.1783 & 0.1667 & 0.2641 & 0.3511 & 0.1954 & 0.3624 & 0.2158 \\
260 & 0.1577 & 0.1569 & 0.3008 & 0.3472 & 0.1644 & 0.2492 & 0.4326 & 0.1817 & 0.2223 & 0.4469 \\
280 & 0.1696 & 0.1549 & 0.1906 & 0.3356 & 0.1711 & 0.2826 & 0.2985 & 0.2777 & 0.2271 & 0.1723 \\
300 & 0.2922 & 0.1563 & 0.1695 & 0.2069 & 0.2684 & 0.3456 & 0.1563 & 0.4015 & 0.1586 & 0.3113 \\
320 & 0.2925 & 0.2251 & 0.1536 & 0.1607 & 0.2973 & 0.4161 & 0.1556 & 0.1960 & 0.1862 & 0.1901 \\
340 & 0.3298 & 0.1775 & 0.5183 & 0.2338 & 0.1831 & 0.1836 & 0.2254 & 0.1711 & 0.2238 & 0.2620 \\
360 & 0.2250 & 0.3286 & 0.3104 & 0.1637 & 0.4682 & 0.2017 & 0.2298 & 0.1599 & 0.1701 & 0.3388 \\
380 & 0.2052 & 0.1531 & 0.2925 & 0.1526 & 0.1847 & 0.2530 & 0.1789 & 0.1797 & 0.3208 & 0.2603 \\
400 & 0.1546 & 0.2975 & 0.1797 & 0.3017 & 0.1595 & 0.2840 & 0.3386 & 0.1510 & 0.4753 & 0.1561 \\
\bottomrule
\end{tabular}
\end{table*}

Three observations, none of them flattering, all of them relevant to how the
reported numbers should be read.

The trajectories are strongly non-monotone. Each seed passes through at least
one evaluation three times worse than its own best, and those collapses occur
at different steps in each run: step 340 for seed 5, step 260 for seed 6, step
360 for seed 7. A single run stopped at a fixed step could land anywhere in
this range, which is why we do not report final-step weights.

Convergence is not established. Seed 6's best evaluation is its last one, so
the 400-step budget may simply be too short; a longer schedule might improve
it further. We report the budget we ran rather than tuning it to the result.

Despite all of this, the three selected checkpoints agree closely when
re-evaluated under one protocol (Table~\ref{tab:main_results}: across-seed
standard deviations between 0.000 and 0.002). The selection procedure is
therefore reproducible even though the path to it is erratic, which is the
strongest claim the evidence supports, and a weaker claim than "training is
stable", which it is not.
